\hypertarget{ia__growfade_8c}{
\section{ia\_\-growfade.c File Reference}
\label{ia__growfade_8c}\index{ia_growfade.c@{ia\_\-growfade.c}}
}
plugin that animates an icon (when starting a program, opening a folder, or somesuch). 

\subsection*{Functions}
\begin{CompactItemize}
\item 
static void \hyperlink{ia__growfade_8c_a0}{ani\_\-icon\_\-fade} (\hyperlink{struct__icon__el}{icon\_\-el} $\ast$el, float alpha, icon\_\-state s)
\begin{CompactList}\small\item\em fade out a single icon. \item\end{CompactList}\item 
void \hyperlink{ia__growfade_8c_a1}{animation} (\hyperlink{struct__icon__list}{icon\_\-list} $\ast$il, \hyperlink{struct__icon}{icon} $\ast$i)
\end{CompactItemize}


\subsection{Detailed Description}
plugin that animates an icon (when starting a program, opening a folder, or somesuch). 

which action is assigned which animation is decided in the theme.

\subsection{Function Documentation}
\hypertarget{ia__growfade_8c_a0}{
\index{ia_growfade.c@{ia\_\-growfade.c}!ani_icon_fade@{ani\_\-icon\_\-fade}}
\index{ani_icon_fade@{ani\_\-icon\_\-fade}!ia_growfade.c@{ia\_\-growfade.c}}
\subsubsection[ani\_\-icon\_\-fade]{\setlength{\rightskip}{0pt plus 5cm}static void ani\_\-icon\_\-fade (\hyperlink{struct__icon__el}{icon\_\-el} $\ast$ {\em el}, float {\em alpha}, icon\_\-state {\em s})\hspace{0.3cm}{\tt  \mbox{[}inline, static\mbox{]}}}}
\label{ia__growfade_8c_a0}


fade out a single icon. 

\begin{Desc}
\item[Parameters:]
\begin{description}
\item[{\em el}]- the icon-element to fade \item[{\em alpha}]- the opaqueness (0.0..1.0) \item[{\em s}]- affect only elements that have one of the states in s\end{description}
\end{Desc}
\hypertarget{ia__growfade_8c_a1}{
\index{ia_growfade.c@{ia\_\-growfade.c}!animation@{animation}}
\index{animation@{animation}!ia_growfade.c@{ia\_\-growfade.c}}
\subsubsection[animation]{\setlength{\rightskip}{0pt plus 5cm}void animation (\hyperlink{struct__icon__list}{icon\_\-list} $\ast$ {\em il}, \hyperlink{struct__icon}{icon} $\ast$ {\em i})}}
\label{ia__growfade_8c_a1}


\begin{Desc}
\item[Parameters:]
\begin{description}
\item[{\em il}]- the list that's affected \item[{\em i}]- the icon we have to highlight\end{description}
\end{Desc}
we can get the evas in question from the icon\_\-list; also, this allows us to mess with icons other than the one selected (fade them out, or somesuch).